% !TeX root = ../main.tex
\section{Analysis}
\label{sec:analysis}

\subsection{Does External Knowledge Need Verification?}
\label{sec:analysis:verification}

RQ3 examines two related questions: whether external knowledge helps under
cross-domain transfer and whether retrieved context should be verified
before it influences the prediction.
We therefore compare the complete model with four controlled train-time
ablations under the same source split, target manifest, and canonical model
seed.

\begin{table*}[t]
    \caption{
    Controlled component ablations on held-out FHM.
    All variants use the canonical seed 42 and are interpreted as diagnostic
    comparisons rather than multi-seed significance estimates.
    }
    \label{tab:ablations}
    \centering
    \renewcommand{\arraystretch}{1.18}
    \begin{tabular*}{0.98\textwidth}{@{\extracolsep{\fill}}lcccccc@{}}
\toprule
Variant & Harm. & Target-P & Target-G & Intent & Tactic & MM-Rel \\
\midrule
% Pending rows must be resolved before the final result freeze.
Ours Full                & -- & -- & -- & -- & -- & -- \\
w/o Retrieval            & 0.473 & 0.228 & 0.099 & 0.057 & -- & 0.174 \\
w/o Support Verifier     & -- & -- & -- & -- & -- & -- \\
w/o Task-Aware Gate      & -- & -- & -- & -- & -- & -- \\
w/o Structured Auxiliary & -- & -- & -- & -- & -- & -- \\
\bottomrule
\end{tabular*}

\end{table*}

The most direct comparison is among the full model,
\textsc{w/o Retrieval}, and \textsc{w/o Support Verifier}.
The first comparison tests whether external context contributes beyond the
internal meme evidence.
The second tests whether simply acquiring external context is sufficient,
or whether task-specific verification provides additional value before
fusion.
The remaining ablations probe whether task-aware routing and structured
auxiliary supervision contribute beyond the external-knowledge pathway.

% TODO_RESULT:
% Interpret Full vs w/o Retrieval vs w/o Support Verifier first.
%
% Case A:
%   Full > w/o Support Verifier > w/o Retrieval
%   -> external context helps, and verification provides additional benefit.
%
% Case B:
%   Full > w/o Retrieval > w/o Support Verifier
%   -> unverified context is sufficiently noisy that using it can be worse
%      than relying on internal evidence alone.
%
% Do not infer multi-seed significance from these canonical-seed ablations.
%
% Final main ablation table target:
% Variant | Harm. | Target-P | Target-G | Intent | Tactic | MM-Rel
% Use Macro-F1 for the six columns; move AUROC and secondary metrics to
% supplementary material.

\subsection{What Does the Verifier Filter?}
\label{sec:analysis:verifier_flow}

RQ4 asks whether Stage~C changes the acquired candidate pool in a meaningful
way rather than acting as a nominal intermediate module.
We analyze the evidence flow from acquisition to verification using the
serialized provenance records.
For each meme, we measure the number of acquired candidates, the number
retained as verified knowledge, the number rejected, and the resulting
acceptance rate.
Where available, we additionally inspect the distribution of candidate
origins and task-specific support states.

These statistics complement the predictive ablation above.
The ablation measures whether verification changes downstream performance,
whereas the candidate-flow analysis characterizes \emph{how} the verifier
changes the information available to the reasoning stage.
Because verifier acceptance is only a model-side screening decision, we do
not interpret it as proof that an accepted candidate is factually correct.

For each meme, we summarize verifier behavior using four primary quantities:
the number of acquired candidates, the number retained as verified evidence,
the number rejected, and the acceptance rate, defined as the fraction of
acquired candidates retained after verification.
We additionally stratify these statistics by candidate origin and, when
available, by task-specific support for target, intent, and rhetorical-tactic
prediction.
All statistics are computed post hoc from the serialized provenance records;
the retrieval corpus, candidate generation procedure, and verification policy
are not modified for this analysis.

% TODO_RESULT:
% Add a compact table or figure only after the audited post-hoc export is
% available. Preferred main-paper statistics:
%
%   acquired candidates / meme
%   verified candidates / meme
%   rejected candidates / meme
%   acceptance rate
%
% Optional if space permits:
%   target-support rate
%   intent-support rate
%   tactic-support rate
%   candidate-origin distribution
%
% Detailed provenance, score distributions, and per-source diagnostics go to
% supplementary material.

\subsection{Structured Effects and Failure Modes}
\label{sec:analysis:structured_errors}
\label{sec:analysis:evidence}

RQ2 and RQ5 concern how the framework behaves beyond binary harmfulness and
where it still fails under domain shift.
We jointly examine the per-task structured results from
Table~\ref{tab:structured-results}, the
\textsc{w/o Task-Aware Gate} and
\textsc{w/o Structured Auxiliary} ablations, and the held-out error cases.
This analysis tests whether components designed for heterogeneous evidence
requirements affect particular structured fields rather than only the
aggregate harmfulness decision.

We also inspect serialized evidence and rationale diagnostics, including
selected-evidence provenance, evidence--rationale overlap, provenance
completeness, and rule-based unsupported-claim indicators.
These quantities are used only as diagnostic properties and are not treated
as direct measurements of causal faithfulness.

Held-out errors are grouped by observable failure mode, including
insufficient contextual knowledge, target confusion, cultural or
event-specific ambiguity, image--text incongruity, rhetorical-tactic
confusion, and irrelevant or insufficiently discriminative retrieved
context.
Representative qualitative examples are included only when they add
information beyond the quantitative error summary; otherwise they are moved
to the supplementary material.

For descriptive error analysis, each inspected case is assigned to the most
directly observable failure category using the model prediction, the
corresponding silver structured annotation, and the serialized evidence
record.
These categories are used to organize recurring failure patterns rather than
as an additional source of ground-truth supervision, and a case is not used
to revise the model, retrieval configuration, or decoding policy.

% TODO_RESULT:
% After result freeze:
% 1. identify structured fields most affected by w/o Task-Aware Gate;
% 2. identify whether structured auxiliary supervision helps the intended
%    fields;
% 3. report the dominant error categories;
% 4. retain only one concise rationale/evidence observation in the main paper
%    unless it materially strengthens the verifier analysis.
